\section{Authoring a problem file}\label{sec:problem-files}

The problem bank is the heart of the system: \emph{one problem per file}, no
preamble, sitting in a topic's \file{problems/} folder. The same file is pulled
into an assessment as a graded question, into a lecture as a worked example, and
into a check-in as a flashcard exercise -- the file itself never changes, only
the import site decides how it appears. Two rules follow from that:

\begin{itemize}
  \item \textbf{A problem file is context-neutral.} It names no label
    (``Q''\,/\,``Example''), no number, and no point value. The consuming class
    supplies the label and number style; the import site supplies the points.
  \item \textbf{A problem file carries what is intrinsically its own}: the
    statement, its hint and solution, and its metadata (topics, tags,
    difficulty, source). Those belong to the problem forever, wherever it is
    used.
\end{itemize}

This intrinsic/extrinsic split is load-bearing. Points are \emph{extrinsic} --
the same problem is worth 5 on one test and 7 on another -- so points are given
where the problem is imported (section~\ref{sec:assessments}), and are
deliberately \emph{not} a metadata field.

\subsection{Anatomy of a problem file}

A complete file, exactly as it sits in the bank:

\begin{manexample}
\ProblemMeta{
  topics     = {differentiation, power-rule},
  tags       = {polynomial, routine},
  difficulty = {basic},
  source     = {own},
}
\begin{problem}
  Differentiate $f(x) = 3x^{4} - 5x^{2} + 2$.
\end{problem}
\begin{hint}
  Differentiate term by term.
\end{hint}
\begin{solution}
  $f'(x) = 12x^{3} - 10x$.
\end{solution}
\end{manexample}

Note what rendered: the metadata block typeset \emph{nothing} (it is
store-only), the problem numbered itself, and -- because this manual forces all
layers on -- the hint and solution both display. In a student build both
disappear without a trace.

To start a new problem, copy an existing file: \file{derivatives/problems/}
in the demo course repo holds one build-tested file of every type.

\subsection{The metadata block}\label{sec:problemmeta}

\cs{ProblemMeta} exists so that ``which problems do I have about the chain
rule?'' is answerable \emph{without running LaTeX} -- nothing has \cs{input}
those files at the moment you are deciding what goes on the test. The
companion script \file{find-problems} answers it by reading the block straight
off disk (section~\ref{sec:find-problems-taste}).

That design puts one discipline on the author -- the block is parsed
\emph{textually}, line by line:

\begin{itemize}
  \item \textbf{One key per line, values braced}, as in the example above. A
    block reflowed onto a single line still \emph{compiles} fine, but the
    problem silently vanishes from the script's index -- the one way to get the
    two readers out of step.
  \item \textbf{Write \texttt{\char`\\\#} for \texttt{\#}} (LaTeX requires it;
    the script strips the backslash back out when reporting, so
    \texttt{source = \{Stewart 3.1 \char`\\\#7\}} reports as
    \texttt{Stewart 3.1 \#7}).
\end{itemize}

The recognised keys are \env{topics}, \env{tags}, \env{difficulty},
\env{source}, and \env{id}. On the LaTeX side they are validated but discarded
-- a mistyped key (\env{topic} for \env{topics}) errors at compile time, so
typos cannot silently un-index a problem. Before inventing a new tag, run
\texttt{find-problems -{}-vocab} to see the vocabulary already in use, and reuse
it.

The block is also optional insurance, not a hard dependency: the exercise
system compiles problem files with or without \file{ProblemMeta.sty} loaded (a
stub swallows the block), so tagging can be adopted -- or abandoned --
without touching the bank.

\subsection{Problems, hints, and solutions}

\env{problem} numbers itself on a document-wide counter and is
cross-referenceable. How the heading \emph{looks} -- ``Problem 7.'' here,
``Q\,7'' with margin marks on a test, a run-in ``Example.'' in lecture notes --
is the consuming class's business, through presentation hooks
(section~\ref{sec:hooks}); the numbering and points bookkeeping stay in the
environment, so no restyle can break the count.

\env{hint} and \env{solution} ride along in the same file, each gated by its
own layer:

\begin{itemize}
  \item \textbf{A solution is the answer key} -- shown only in a solutions
    build, set in the green box you saw above. \env{solution}\texttt{[notitle]}
    drops the corner tab, leaving a bare box.
  \item \textbf{A hint is student-facing help} -- shown only in a hints build,
    as a light run-in (deliberately plainer than a solution: a nudge, not the
    answer). It sits inline where written, not collected at the end.
\end{itemize}

The two layers are independent: a hints build is not a solutions build, and a
per-import key can force one problem's hint on or off against the document
toggle (section~\ref{sec:assessments}).

\subsection{Dependent parts in one file: \env{qparts}}

When parts genuinely depend on each other -- one stem, sub-questions (a), (b),
(c) -- they are \emph{one} problem and belong in \emph{one} file, with the
parts marked by \cs{qpart}:

\begin{manexample}
\begin{problem}
  Evaluate.
  \begin{qparts}
    \qpart Find $\lzoo{x}{x^{2}e^{x}}$.
    \begin{solution}[notitle]
      $2x\,e^{x} + x^{2}e^{x} = x\,e^{x}(x+2)$.
    \end{solution}
    \qpart Find $\lzoo{x}{x^{3}e^{x}}$.
  \end{qparts}
\end{problem}
\end{manexample}

Each part letters itself, and a solution (or hint) can follow its own
\cs{qpart}. Per-part points and per-part workspace are supplied at the import
site, for example \texttt{points=\{2,1\}} and \texttt{workspaces=\{4,4\}} --
keeping even part-level grading extrinsic (section~\ref{sec:assessments}).

The other thing people call ``parts'' -- \emph{independent} problems grouped
under a shared instruction (``Evaluate the following:'') -- is not this. Those
stay atomic one-problem files and are grouped at assembly time with
\env{problemgroup} (section~\ref{sec:assessments}).

\subsection{Multiple choice}

An MC question is \emph{not} a separate problem type: it is a normal
\env{problem} whose body carries a \env{choices} list -- so it numbers, grades,
tags and imports exactly like everything else. Mark the key with the optional
argument:

\begin{manexample}
\begin{problem}
  What is the domain of $f(x)=\dfrac{\sqrt{x+1}}{x-7}$?
  \begin{choices}
    \choice{$\{x : x \ge -1\}$}
    \choice{$\{x : x > -1\}$}
    \choice[correct]{$\{x : x \ge -1,\ x \ne 7\}$}
    \choice{$\{x : x > -1,\ x \ne 7\}$}
  \end{choices}
\end{problem}
\end{manexample}

The key is boxed \emph{only} in a solutions build (the green cue above -- gated
by the same toggle as \env{solution}, so the answer never leaks into a student
copy), and swapping which option is correct is an in-place edit. Marking
several options \texttt{[correct]} makes a select-all-that-apply.

One environment, three layouts:

\begin{manexample}
Is $f(x) = x^{3}$ even, odd, both, or neither?
\begin{choices}[onepar]
  \choice{Even}
  \choice[correct]{Odd}
  \choice{Both}
  \choice{Neither}
\end{choices}
\end{manexample}

\begin{itemize}
  \item \emph{(default)} -- one option per line, for long options;
  \item \texttt{[onepar]} -- inline, for short options as above;
  \item \texttt{[grid]} -- a grid, \cs{ChoicesGridCols} wide (default 2),
    row-major.
\end{itemize}

For a grid of \emph{graphs}, \cs{graphchoice}\texttt{[correct]\{FILE\}} is a
\cs{choice} that sizes its image to the column and sets the letter beside it.
This manual has no image files, so the grid below is source-only; the live
version is \file{mc-graph-grid.tex} in the demo repo:

\begin{manexsource}
\begin{problem}
  The graph of $y = f(x)$ has a horizontal tangent at $x = 1$. Which of
  the following could be the graph of $y = f'(x)$?
  \begin{choices}[grid]
    \graphchoice{../images/cubic}
    \graphchoice[correct]{../images/tangent}
    \graphchoice{../images/parabola}
    \graphchoice{../images/step}
  \end{choices}
\end{problem}
\end{manexsource}

\subsection{Matching}

A matching question is likewise a normal \env{problem}. Left items carry the
answer key intrinsically -- each names its match in the optional argument, shown
(green-boxed) only in a solutions build:

\begin{manexample}
Match each function with its derivative.
\begin{matching}
  \leftitem[B]{$\sin x$}
  \leftitem[C]{$\cos x$}
  \leftitem[A]{$\tan x$}
  \rightcolumn
  \rightitem{$\sec^{2} x$}
  \rightitem{$\cos x$}
  \rightitem{$-\sin x$}
  \rightitem{$\sec x \tan x$}
\end{matching}
\end{manexample}

Left items letter in lowercase roman, right items in uppercase alphabetic --
two visibly different systems, so ``match [i] with [C]'' can never be
misread. The columns may have different lengths (distractors on the right, as
above), and \cs{rightcolumn} splits them with a wide gap for students to draw
their connecting lines in. Graph items reuse the \cs{graphchoice} sizing:
\cs{leftgraph}\texttt{[X]\{FILE\}} / \cs{rightgraph}\texttt{\{FILE\}} put
graphs in either column. The demo's \file{match-graphs-phrases.tex} keeps the
graphs on the \emph{right} so the answer key rides the roomy text column rather
than a column-width image; \file{match-trig-derivatives.tex} matches text to
text. There are deliberately no drawn answer arrows -- the inline key on the
left item is the answer, and skipping arrows keeps the system TikZ-free.

\subsection{Finding problems in the bank}\label{sec:find-problems-taste}

The payoff for the metadata discipline -- from anywhere in the course repo:

\begin{manexsource}
$ find-problems --topic chain-rule --difficulty basic
$ find-problems --vocab      # the tag vocabulary in use, with counts
$ find-problems --untagged   # the tagging backlog
$ find-problems --import     # print matches as paste-ready \import lines
\end{manexsource}

Filters combine (AND), match case-insensitively as regular expressions, and
never run LaTeX. The full command reference is in
section~\ref{sec:build-and-find}.
